\documentclass[11pt]{article}
\usepackage[margin=1in]{geometry}
\usepackage{amsmath,amssymb}
\usepackage{booktabs}
\usepackage{longtable}
\usepackage{array}
\usepackage{siunitx}
\usepackage{hyperref}
\usepackage{xcolor}
\usepackage{enumitem}
\hypersetup{colorlinks=true,linkcolor=blue,urlcolor=blue}

\newcommand{\chizxx}{\chi^{(0)}_{z;xx}}
\newcommand{\lamR}{\lambda_{R}}
\newcommand{\muB}{\mu_{B}}

\title{Replication Report:\\
\emph{Geometry-Driven Nonlinear Orbital Magnetoelectric Effect}\\
\large Jia, Qiao, Wang --- arXiv:2605.17462}
\author{Replication study (\texttt{jia2026}) --- REPLICATE-PROJECT / TEXTURE-orbital}
\date{2026-07-17}

\begin{document}
\maketitle

\begin{abstract}
We attempt a numerical reproduction of the intrinsic nonlinear orbital
magnetoelectric effect (NOME) reported by Jia \emph{et al.} for their
modified Kane--Mele honeycomb model (Model~1, Eq.~11). We build a
from-scratch \texttt{numpy} tight-binding implementation of the 4-band
Bloch Hamiltonian and evaluate the conventional matrix-element
(``od'') contribution to $\chizxx$ as a Brillouin-zone integral, together
with the analogous spin response. \textbf{Verdict: PARTIAL.} We cleanly
reproduce (a) the 4-band structure of Fig.~1(a) ($\sim$5~eV bandwidth,
$\sim$18~meV central gap); (b) the flagship \emph{symmetry} statement that
$\chizxx$ is an \emph{even} function of $\lamR$ (P-even) to machine
precision ($\sim\!10^{-16}$), with a nonzero value at $\lamR=0$; and (c)
the qualitative orbital-vs-spin \emph{opposite-sign} relationship near the
gap. We do \textbf{not} reproduce the headline quantitative claim that the
orbital moment is $\approx 3\times$ the spin moment, nor the absolute
$\muB/\mathrm{V}^2$ prefactor: both depend on the geometric /
Hermitian-connection response sector (Eqs.~4--9), the paper's actual novel
physics, which was not implemented in this compute pass. This is an honest
partial reproduction that \emph{validates the symmetry structure} of the
effect but not its flagship magnitude.
\end{abstract}

\section{Paper summary}
The paper proposes a \emph{nonlinear} orbital magnetoelectric effect
(NOME): a magnetization response second order in the electric field,
$\delta M^{c} = \chi^{c;ab} E_a E_b$, described by a rank-3 P-even
pseudotensor. Unlike the \emph{linear} OME, the NOME is symmetry-allowed
even in centrosymmetric (inversion-symmetric) materials, dramatically
enlarging the set of hosting magnetic point groups (MPGs): 41 MPGs support
the intrinsic NOME and 55 the extrinsic NOME in 2D, versus only 10/8 for
the linear OME. The intrinsic response decomposes into three geometric
contributions --- a conventional matrix-element (``od'') term, a dipole
(``d'') term, and a positional-shift (``ic'') term --- all governed by the
quantum-geometric \emph{Hermitian connection}. The authors argue the
geometric $d+ic$ terms \emph{dominate} near the band edge. Two model
demonstrations are given: (Model~1) a modified Kane--Mele honeycomb
(MPG $6m'm'$) for the intrinsic effect, predicting orbital
$M \approx 3\times$ spin $M$ with opposite sign and a nonzero
$\sim\!10\,\muB/\mathrm{V}^2$ minimum at $\lamR=0$; and (Model~2) a CuMnAs
antiferromagnet ($PT$-symmetric) for the extrinsic effect. Estimated
magnitudes ($\sim\!10^{-5}\,\muB/\mathrm{nm}^2$ for the honeycomb at
$E=10^5$~V/m) are claimed detectable by polar MOKE.

\section{Claims tested}
Table~\ref{tab:claims} lists the paper's central claims, whether each is
testable with our compute, whether we tested it, and the paper vs.\
reproduced values.

\begin{longtable}{@{}p{1.6cm} p{4.3cm} c c p{2.6cm} p{2.9cm}@{}}
\caption{Claims, testability, and reproduction outcome.}
\label{tab:claims}\\
\toprule
\textbf{Claim} & \textbf{Statement} & \textbf{Testable} & \textbf{Tested}
& \textbf{Paper value} & \textbf{Reproduced} \\
\midrule
\endfirsthead
\toprule
\textbf{Claim} & \textbf{Statement} & \textbf{Testable} & \textbf{Tested}
& \textbf{Paper value} & \textbf{Reproduced} \\
\midrule
\endhead
C4-bands & KM honeycomb 4-band structure, meV gaps (Fig.~1a) & yes & yes
& 4 bands, meV gaps near K/K$'$, $\sim$eV width
& \textbf{MATCH}: 4 bands, $E\in[-2560,2560]$~meV, gap 18.04~meV \\
\addlinespace
C4-even-$\lamR$ & $\chizxx$ even in $\lamR$ (P-even), nonzero min at
$\lamR{=}0$ & yes & yes & even in $\lamR$; $\sim\!10\,\muB/\mathrm{V}^2$
at $\lamR{=}0$ & \textbf{MATCH}: even to rel.\ err $9.1\times10^{-16}$;
nonzero raw value $13.94$ at $\lamR{=}0$ \\
\addlinespace
C4-orb/spin & orbital $\approx 3\times$ spin, opposite sign & yes & yes
(od only) & orbital $\approx 3\times$ spin, opposite sign
& \textbf{PARTIAL}: sign opposite near gap \checkmark; ratio (od-only)
median $\approx 11.6$, peak $-115.6$ (\emph{not} $3\times$) \\
\addlinespace
C4-onlyxx & $6m'm'$: only $\chi_{z;xx}{=}\chi_{z;yy}$ finite & yes & no
& $\chi_{z;xy}{=}0$ by symmetry
& \textbf{Not enumerated}; only allowed $xx$ computed \\
\addlinespace
C6-order & magnitude $\sim\!10\,\muB/\mathrm{V}^2$
($\sim\!10^{-5}\muB/\mathrm{nm}^2$) & yes & partial & $\sim\!10\,\muB/\mathrm{V}^2$
& \textbf{Not pinned}: structure reproduced, absolute prefactor not fixed \\
\addlinespace
C1/C2/C3/C5 & existence, 3-term geometric split, MPG count, CuMnAs model
& yes & no & --- & \textbf{Not attempted} this pass \\
\bottomrule
\end{longtable}

\section{Method}
\paragraph{Model.} We implement Model~1, the modified Kane--Mele
honeycomb (Eq.~11), on a 4-state basis (2 sublattice $\times$ 2 spin):
\[
H = -t\sum_{\langle ij\rangle} c^\dagger_i c_j
   + i\lamR \sum_{\langle ij\rangle} (\boldsymbol{\sigma}\times \mathbf{d}_{ij})_z
   + i\lambda_{so}\sum_{\langle\langle ij\rangle\rangle} \nu_{ij}\,\sigma_z
   + \lambda \sum_i s_z ,
\]
comprising nearest-neighbor hopping, inversion-breaking Rashba SOC,
Kane--Mele mirror SOC (opposite sign on the two sublattices), and a
$T$-breaking exchange term. Parameters follow the Fig.~1 caption:
$t = 0.85$~eV, $\lamR = 20$~meV, $\lambda = 10$~meV,
$\lambda_{so} = 10$~meV, $T = 20$~K, evaluated on a $120\times120$
Brillouin-zone mesh.

\paragraph{Response.} From the Bloch Hamiltonian we compute, per
$k$-point: eigenenergies; velocity matrices $v^a = U^\dagger(\partial_a H)U$
(central finite difference in $k$); the interband Berry connection
$r^a_{nm} = i\,v^a_{nm}/(E_n - E_m)$; the orbital-moment matrix
$L^z_{nm} = \tfrac14\varepsilon_{zab}\{r^a,v^b\}_{nm}$; the quantum metric
$G^{xx}_{nm}$; and the spin matrix $S^z_{nm}$. We then evaluate the
conventional intrinsic (``od'') NOME term $\chi^{(0,od)}_{z;xx}$ --- the
first line of Eq.~(3) --- as a Fermi--Dirac-weighted BZ integral, and the
analogous \emph{spin} response $\chi^{s(0,od)}_{z;xx}$ via the paper's
stated substitution $L^c \to s^z$. Sweeps: band structure along
K$'$--$\Gamma$--M--K--$\Gamma$ (Fig.~1a); $\chi^{(0,od)}$ vs.\ $\mu$
(Fig.~1b); and $\chi^{(0,od)}$ vs.\ $\lamR$ at fixed $\mu=50$~meV
(Fig.~1c). All logic is in \texttt{work/reproduce.py}; numeric outputs are
serialized to \texttt{work/results.json} and figures to
\texttt{work/figs/*.png}.

\paragraph{What we deliberately did not implement.} The geometric ``d''
and ``ic'' terms (Eqs.~4--6, 8, 9), which require the full Hermitian
connection $C^{\alpha ab}_{nm}$, positional-shift / covariant-derivative
tensors ($N^{\alpha ab}_{nm}$, $L^{\alpha ab}_{nm}$, $D^a_{mn}$), were out
of the compute budget for a defensibly gauge-invariant implementation. As
discussed below, this is precisely the sector responsible for the flagship
$3\times$ magnitude.

\section{Results}

\subsection{Band structure (C4-bands): MATCH}
The 4-band structure reproduces Fig.~1(a): four bands spanning
$E\in[-2560,2560]$~meV ($\approx 6t = 5.1$~eV total bandwidth), with a
central (band 2$\leftrightarrow$3) gap of $18.04$~meV set by the combined
SOC + exchange scales. See \texttt{work/figs/fig1a\_bands.png}.

\subsection{P-even / even-in-$\lamR$ symmetry (C4-even-$\lamR$): MATCH (machine precision)}
This is the strongest reproduced result. We find
\[
\chi^{(0,od)}_{z;xx}(+\lamR) = \chi^{(0,od)}_{z;xx}(-\lamR)
\quad\text{to relative error } 9.08\times10^{-16},
\]
verified across $\lamR = 10, 20, 30$~meV (Table~\ref{tab:even}). The
response is \emph{nonzero at $\lamR = 0$} (raw value $13.94$),
structurally reproducing the paper's statement of a nonzero minimum when
inversion is restored. This directly and cleanly confirms the P-even
character of the intrinsic NOME. See \texttt{work/figs/fig1c\_chi\_vs\_lamR.png}.

\begin{table}[h]
\centering
\caption{Even-in-$\lamR$ check: $\chi^{(0,od)}_{z;xx}(+\lamR)$ vs.\
$\chi^{(0,od)}_{z;xx}(-\lamR)$ (raw units).}
\label{tab:even}
\begin{tabular}{@{}c c c@{}}
\toprule
$\lamR$ (eV) & $\chi(+\lamR)$ & $\chi(-\lamR)$ \\
\midrule
0.01 & $+2.193426222912495$  & $+2.1934262229124903$ \\
0.02 & $-5.2896687883299345$ & $-5.289668788329926$  \\
0.03 & $-0.2066981391265077$ & $-0.20669813912650806$ \\
\bottomrule
\end{tabular}
\end{table}

\subsection{Orbital vs.\ spin (C4-orb/spin): PARTIAL}
Near the band-gap region the orbital ``od'' response has the
\emph{opposite sign} to the spin response, consistent with the paper's
qualitative claim. However, the \emph{quantitative} magnitude is not
reproduced: the od-only orbital/spin ratio is
\[
\text{median (near cond.\ edge)} \approx 11.6,\qquad
\text{peak} \approx -115.6,
\]
i.e.\ $\sim\!10$--$100\times$, rather than the claimed $\approx 3\times$.
As explained in the critique, the $3\times$ figure is a property of the
\emph{total} (od $+$ geometric) response; the od term alone overshoots.
See \texttt{work/figs/fig1b\_chi\_vs\_mu.png}.

\subsection{Absolute magnitude (C6): not pinned}
We report the response in raw/internal geometric units. Converting to
$\muB/\mathrm{V}^2$ (and thence to $\sim\!10^{-5}\,\muB/\mathrm{nm}^2$ at
$E=10^5$~V/m) requires the full dimensional prefactor bookkeeping of
Eq.~(1) (electron charge, lattice length, $\hbar$, cell area, the E-field
length convention). We reproduce the response \emph{structure} but do not
fix the absolute scale; C6 is neither confirmed nor refuted.

\section{Per-claim: what worked / what did not}
\begin{itemize}[leftmargin=1.4em]
\item \textbf{C4-bands --- worked.} 4 bands, correct bandwidth, meV gap.
\item \textbf{C4-even-$\lamR$ --- worked (best result).} Even symmetry to
$\sim\!10^{-16}$; nonzero at $\lamR=0$. This validates the core symmetry
structure of the effect.
\item \textbf{C4-orb/spin --- partial.} Sign relationship reproduced;
$3\times$ magnitude \emph{not} reproduced (od-only ratio $\sim\!10$--$100\times$).
\item \textbf{C4-onlyxx --- not done.} Only the allowed $xx$ component was
computed; forbidden-component vanishing was not enumerated.
\item \textbf{C6 --- not pinned.} Absolute $\muB/\mathrm{V}^2$ prefactor
not fixed (dimensional bookkeeping).
\item \textbf{C5 (CuMnAs Model~2) --- not attempted.}
\item \textbf{C3 (MPG enumeration, 41/55) --- not attempted.}
\end{itemize}

\section{Honest critique}
The strongest reproduced result is the \emph{symmetry} statement: that the
intrinsic NOME $\chizxx$ is an even function of $\lamR$ (P-even), which we
confirm to machine precision, together with a nonzero value at $\lamR=0$.
That is a genuine, clean validation of the effect's symmetry backbone.

However, the paper's \emph{flagship} quantitative result --- orbital
moment $\approx 3\times$ spin moment --- was \textbf{not} quantitatively
reproduced. The reason is specific and important: the $3\times$ ratio is a
property of the \emph{total} response, which is dominated by the geometric
$d+ic$ / Hermitian-connection sector (Eqs.~4--9). That sector \emph{is}
the paper's actual novel physics, and it was \emph{not coded} in this
compute pass. Our od-only implementation --- the one robustly
well-defined, gauge-invariant term --- overshoots the ratio to
$\sim\!10$--$100\times$ precisely because it omits the geometric
contributions that would rebalance orbital vs.\ spin near the edge. In
short: \textbf{this replication validates the symmetry structure of the
NOME but not its flagship magnitude}, and the missing piece is exactly the
paper's headline geometric mechanism. Any claim of full reproduction would
be dishonest without implementing Eqs.~(4)--(6) with correct covariant
derivatives and a Dirac-point convergence study.

Two further honesty caveats: (i) the absolute $\muB/\mathrm{V}^2$
prefactor (C6) is not pinned, so the ``$\sim\!10\,\muB/\mathrm{V}^2$''
order-of-magnitude is unverified; and (ii) the $\chi$-vs-$\lamR$ curve at
fixed $\mu$ is mesh-sensitive/noisy (the Fermi level crosses band edges as
$\lamR$ varies on a fixed $120\times120$ grid) --- the \emph{even
symmetry} is exact regardless, but the fine shape would need a denser,
adaptively-refined mesh near K/K$'$.

\section*{Open Questions}
\addcontentsline{toc}{section}{Open Questions}
\begin{description}[leftmargin=2.2em,style=nextline]
\item[Q1.] Does implementing the full Hermitian-connection geometric
sector (Eqs.~4--6, 8, 9) recover the paper's $\approx 3\times$ orbital/spin
ratio? Our od-only ratio is $\sim\!10$--$100\times$; the missing geometric
$d+ic$ terms are the hypothesized rebalancing mechanism.
\item[Q2.] Is the machine-precision $\lamR$-evenness of $\chizxx$ a
property of the \emph{total} response, or only of the od sector we
computed? If the geometric terms break or preserve the evenness at the
same precision, that is a strong internal-consistency test of the P-even
claim.
\item[Q3.] What fixes the absolute $\muB/\mathrm{V}^2$ prefactor (C6)?
Which dimensional convention in Eq.~(1) (E-field length gauge, cell area,
$\hbar$/charge factors) reconciles our raw units with the claimed
$\sim\!10\,\muB/\mathrm{V}^2$ and $\sim\!10^{-5}\,\muB/\mathrm{nm}^2$?
\item[Q4.] How does the extrinsic CuMnAs Model~2 (Eq.~12, $PT$-symmetric)
compare --- does its single allowed $\chi^{(1)}_{z;xy}$, dominated by the
Hermitian connection with a Dirac-point peak, reproduce cleanly where the
intrinsic $3\times$ did not?
\item[Q5.] How sensitive is the $3\times$ ratio (and the whole
$\chi(\mu)$ profile) to temperature $T$ and chemical potential $\mu$, and
to BZ-mesh density near the Dirac points? Does an adaptive $\geq200\times200$
mesh remove the observed $\lamR$-sweep noise and stabilize the ratio?
\end{description}

\section{Verdict}
\textbf{PARTIAL.} We reproduced the model, its band structure, and --- to
machine precision --- the paper's central \emph{symmetry} statement that
the intrinsic NOME $\chizxx$ is an even function of $\lamR$ (P-even) with a
nonzero value at $\lamR=0$, plus the qualitative orbital-vs-spin
opposite-sign relationship. We did \emph{not} reproduce the flagship
quantitative claim (orbital $\approx 3\times$ spin) or the absolute
$\muB/\mathrm{V}^2$ magnitude, because both depend on the geometric /
Hermitian-connection sector (Eqs.~4--9) --- the paper's genuine novel
physics --- which was deliberately not implemented in the compute budget,
plus dimensional-prefactor bookkeeping. Model~2 (CuMnAs) and the MPG
enumeration were not attempted. This is an honest partial that validates
the effect's symmetry structure but not its flagship magnitude.


\section{Verdict}
\textbf{Verdict: PARTIAL} (Coverage 5/10, Agreement 5/10). — 3-judge consensus (median cov/agr, majority verdict): Replication covers Model-1 band structure and verifies χz;xx is λR-even with nonzero value at λR=0, but it does not reproduce the paper’s key quantitative magnitude claims (orbital≈3×spin and ~10 μB/V^2 scale) and omits the dominant geometric/Hermitian-connection terms plus other core results (tensor-component symmetry check, MPG counts, Model-2).
% census-verdict: PARTIAL assigned 2026-07-18 by LLM judge (Argo Opus)

\end{document}
